\documentclass[UTF8,a4paper,11pt,fontset=fandol]{ctexart}
\usepackage[margin=23mm,headheight=15pt]{geometry}
\usepackage{amsmath,amssymb,booktabs,longtable,array,calc,graphicx,xcolor}
\usepackage{fancyhdr,fvextra,upquote,xurl,caption,etoolbox,needspace,float}
\newcounter{none}
\usepackage[unicode,colorlinks=true,linkcolor=ReportBlue,urlcolor=ReportBlue]{hyperref}
\definecolor{ReportBlue}{HTML}{24516A}
\setmainfont{texgyretermes-regular.otf}[BoldFont=texgyretermes-regular.otf]
\setsansfont{texgyreheros-regular.otf}[BoldFont=texgyreheros-regular.otf]
\setmonofont{lmmono10-regular.otf}[Scale=0.85,BoldFont=lmmono10-regular.otf]
\setCJKmainfont{FandolSong-Regular.otf}[BoldFont=FandolSong-Regular.otf]
\setCJKsansfont{FandolHei-Regular.otf}[BoldFont=FandolHei-Regular.otf]
\ctexset{section={format=\Large\normalfont\color{ReportBlue}},
 subsection={format=\large\normalfont\color{ReportBlue}},
 subsubsection={format=\normalsize\normalfont}}
\setlength{\parindent}{2em}
\setlength{\parskip}{4pt}
\linespread{1.16}
\setlength{\emergencystretch}{3em}
\widowpenalty=10000
\clubpenalty=10000
\setlength{\LTpre}{6pt}
\setlength{\LTpost}{6pt}
\AtBeginEnvironment{longtable}{\small\renewcommand{\arraystretch}{1.2}}
\DefineVerbatimEnvironment{verbatim}{Verbatim}{breaklines,breakanywhere,fontsize=\footnotesize}
\providecommand{\tightlist}{\setlength{\itemsep}{2pt}\setlength{\parskip}{0pt}}
\providecommand{\pandocbounded}[1]{#1}
\makeatletter
\def\fps@figure{H}
\def\maxwidth{\ifdim\Gin@nat@width>\linewidth\linewidth\else\Gin@nat@width\fi}
\def\maxheight{\ifdim\Gin@nat@height>0.72\textheight0.72\textheight\else\Gin@nat@height\fi}
\makeatother
\setkeys{Gin}{width=\maxwidth,height=\maxheight,keepaspectratio}
\captionsetup{font=small,labelfont=normalfont,margin=8mm}
\pagestyle{fancy}
\fancyhf{}
\fancyfoot[C]{\small\thepage}
\renewcommand{\headrulewidth}{0pt}
\fancypagestyle{plain}{
 \fancyhf{}
 \fancyfoot[C]{\small\thepage}
 \renewcommand{\headrulewidth}{0pt}
}
\hypersetup{pdftitle={基于机器学习的网络流量入侵检测系统的设计与实现},
 pdfauthor={},pdfsubject={课程实验报告：离线模型比较与持续学习模拟},pdfkeywords={UNSW-NB15,MLP,Transformer}}
\begin{document}
\begin{titlepage}
\vspace*{18mm}
{\large\color{ReportBlue}课程设计 · 最终实验报告\par}
\vspace{24mm}
{\fontsize{25}{38}\selectfont 基于机器学习的网络流量\par 入侵检测系统的设计与实现\par}
\vspace{14mm}
{\large 模型结构与实验比较\quad /\quad 实时推理与定时训练模拟\par}
\vspace{5mm}
{\large Web Dashboard 与可复现证据\par}
\vfill
\noindent\rule{\textwidth}{0.5pt}
\vspace{5mm}
\noindent 数据集：UNSW-NB15\par
\noindent 实验：2026 年 11 月 3 日至 6 日\par
\noindent 整理与验证：2026 年 11 月 13 日\par
\noindent 公开版本：不含个人姓名、学号或演示文稿\par
\vspace{12mm}
\end{titlepage}
\pagenumbering{roman}
\begingroup
\small\linespread{1}\selectfont
\setlength{\parskip}{0pt}
\tableofcontents
\endgroup
\clearpage
\pagenumbering{arabic}
\section{摘要}\label{ux6458ux8981}

本项目面向网络流量正常/攻击二分类，完成了从公开数据获取、特征审计、训练与验证、模型保存到独立推理的实验系统，并实现由
CLI 启动、长期运行的本地 Web Dashboard。系统使用 UNSW-NB15
的预定义训练/测试角色，在隔离冲突标签和跨集合重复特征后，以 131,706
条记录训练、32,927 条记录验证，在固定的 82,332
条测试记录上比较逻辑回归、随机森林、Torch 线性模型、MLP 与轻量表格
Transformer。

在已完成的模型族评估中，轻量 Transformer 的测试 Accuracy 为
88.69\%、攻击类 F1 为 90.54\%、误报率为 23.06\%。随机森林与原 MLP 的 F1
分别为 89.20\% 和 89.16\%。调参 MLP 的验证 F1 略升，但测试 F1 降至
87.49\%，说明验证集上的微小提升不能直接等同于泛化改善。报告保留调参前后结果，不以测试成绩重新选择候选；2026-11-05
调参实验的整体验证获选模型仍为 MLP。

实时部分以真实开发集记录模拟特征流到达：常驻模型先预测，延迟揭示标签，按来源去重入池，定时冻结累计样本并在独立进程重训，通过固定验证门槛后切换模型。历史演示包含
20,412 次检测事件和 8
次训练任务，验证了候选接受与拒绝路径。回放包含重复事件，其指标不作为独立测试效果。系统适用于课程实验和机制展示，尚无真实网络采集或生产部署证据。

关键词：网络入侵检测；UNSW-NB15；MLP；表格
Transformer；持续学习模拟；模型验证。

阅读入口：\href{https://github.com/ZureTz/ml-network-intrusion-detection/blob/master/reports/final/report.pdf}{最终
PDF}、\href{https://github.com/ZureTz/ml-network-intrusion-detection/blob/master/reports/final/report.tex}{独立
LaTeX
源文件}、\href{https://github.com/ZureTz/ml-network-intrusion-detection/blob/master/reports/final/README.md}{报告构建说明}。本报告为公开版，不含个人姓名、学号和
PPT。实验主要完成于 2026 年 11 月 3 日至 6 日，本版整理与工程验证日期为
2026 年 11 月 13 日。

\section{问题定义与系统需求}\label{ux95eeux9898ux5b9aux4e49ux4e0eux7cfbux7edfux9700ux6c42}

每条输入是一行已提取的流量特征，而非原始报文或任意系统日志。输入包含 39
个数值字段和 3
个类别字段；输出为正常/攻击标签以及攻击概率，标签映射为正常 0、攻击
1，攻击为正类。预测概率是模型输出，尚未进行概率校准。

系统的离线流程为
\texttt{prepare\ →\ train\ →\ evaluate\ /\ predict}，在线形式的模拟入口为
\texttt{serve}。需求包括：预处理仅拟合训练集；候选仅由验证集选择；模型能在新进程独立加载；无标签输入可分批预测；异常输入给出明确错误；运行状态、训练记录及模型版本可追溯。模拟模式还要求模型常驻内存、训练不阻塞检测、延迟标签不提前用于学习、重启后能恢复实验。

当前实现不执行实时抓包，也不从报文计算特征。模拟器重放已有真实记录，不生成新的独立网络样本。定时训练是在累计已确认样本上从头重新拟合模型，不是对每条到达记录立即执行增量梯度更新。合成数据仅用于自动化测试，未用于正式检测效果结论。

\section{数据集、特征与防泄漏处理}\label{ux6570ux636eux96c6ux7279ux5f81ux4e0eux9632ux6cc4ux6f0fux5904ux7406}

\subsection{来源与字段}\label{ux6765ux6e90ux4e0eux5b57ux6bb5}

UNSW 官方介绍的预定义训练集为 175,341 行，测试集为 82,332
行，包含正常流量及九类攻击背景；本实验把攻击统一映射为正类，未实现九类攻击的多分类识别。原始全量数据的
49 字段说明与预定义划分 CSV 的 45 列布局不同，本项目使用后者的 42
个输入字段。\href{https://research.unsw.edu.au/projects/unsw-nb15-dataset}{官方数据页}提供数据说明与引用要求。

2026-11-03 获取时官方下载入口要求登录，实际使用固定版本的
\href{https://huggingface.co/datasets/Mireu-Lab/UNSW-NB15/tree/3eb02c4a0a29866b7abcb5ef77c45cf4fcc8f6b0}{Mireu-Lab
镜像}。该镜像的 \texttt{test.csv} 对应本项目训练角色，\texttt{train.csv}
对应测试角色；映射依据行数、类别分布与镜像间哈希交叉核验。原始下载字节不变，只重命名文件；尚未与官方原始文件完成逐字节核验。完整
SHA-256、字段含义、授权范围及下载方式见
\href{https://github.com/ZureTz/ml-network-intrusion-detection/blob/master/reports/dataset.md}{数据说明}，本仓库不分发完整数据集。

\texttt{proto}、\texttt{service}、\texttt{state}
为类别字段。数值字段包括持续时间、包数、字节数、速率、TTL、丢包、窗口、时延、抖动及邻近连接统计。\texttt{id}、\texttt{label}、\texttt{attack\_cat}
不作为特征；\texttt{source\_file}、\texttt{source\_row}
只用于追溯。攻击类别只用于评估后的错误细分，不能从中派生训练输入。

\subsection{审计与划分}\label{ux5ba1ux8ba1ux4e0eux5212ux5206}

原始训练角色中，同特征冲突标签涉及 229 组、940 行；测试角色有 6 组、14
行。默认严格策略会拒绝冲突。正式配置显式采用
\texttt{quarantine\_development}，隔离训练侧冲突记录，完整保留测试侧冲突记录，并保留来源清单。随后移除开发集中与测试集特征完全相同的
9,768 行，得到 164,633 行开发数据。

使用固定字段顺序和规范类型计算特征指纹，以种子 42 的
\nolinkurl{StratifiedGroupKFold(n_splits=5)}
第一折作为验证集。同一特征组不跨训练/验证；同标签的组内重复记录仍保留，不能声称重复样本影响已全部消除。

\Needspace{6\baselineskip}

{\def\LTcaptype{none} % do not increment counter
\begin{longtable}[]{@{}lrrrl@{}}
\toprule\noalign{}
集合 & 总行数 & 正常 & 攻击 & 用途 \\
\midrule\noalign{}
\endhead
\bottomrule\noalign{}
\endlastfoot
训练 & 131706 & 43237 & 88469 & 拟合预处理和分类模型 \\
验证 & 32927 & 10809 & 22118 & 选择候选，不参与拟合 \\
测试 & 82332 & 37000 & 45332 & 冻结候选后的效果评估 \\
\end{longtable}
}

三个实验批次的划分文件 SHA-256 相同，见
\href{https://github.com/ZureTz/ml-network-intrusion-detection/blob/master/reports/results/full-20261103/data/manifest.json}{历史清单}、\href{https://github.com/ZureTz/ml-network-intrusion-detection/blob/master/reports/results/torch-full-20261105/data/manifest.json}{Torch
清单}和\href{https://github.com/ZureTz/ml-network-intrusion-detection/blob/master/reports/results/torch-tune-20261105/data/manifest.json}{调参清单}。因此结果可在这一固定基准上比较，但不构成不同训练方法的等预算因果实验。

\subsection{预处理与有效性边界}\label{ux9884ux5904ux7406ux4e0eux6709ux6548ux6027ux8fb9ux754c}

Torch 模型使用
\texttt{ColumnTransformer}：数值字段以训练集中位数填充并标准化，类别字段
One-Hot
编码；未知类别允许编码但记录数量。变换器只在训练集拟合，并与分类器共同构成
Pipeline。完整训练集拟合后输入维度为 193，其中数值 39 维、类别编码 154
维。历史逻辑回归同样标准化，随机森林不依赖数值标准化；历史实现以其原始配置和代码为准。

CSV 缺少可靠时间、原始
IP、主机和会话键，不能证明这些层面的隔离。TTL、TCP
序列号和连接聚合统计可能反映实验室采集环境，聚合窗口无法由当前 CSV
重建。测试特征曾用于重叠审计，虽未用于拟合和选参，仍有别于完全不接触测试输入的协议。指纹分组还存在极低概率的哈希碰撞风险。

\section{系统设计与模型结构}\label{ux7cfbux7edfux8bbeux8ba1ux4e0eux6a21ux578bux7ed3ux6784}

\subsection{模块与数据流}\label{ux6a21ux5757ux4e0eux6570ux636eux6d41}

数据模块负责字段契约、审计和划分；训练模块负责 Pipeline
拟合、验证选择、评估及持久化；推理模块负责可信模型加载、输入校验和分批预测；模拟引擎负责到达调度、反馈标签、训练进程、状态持久化和模型切换；FastAPI
与前端负责展示和控制。

离线实验按候选逐个训练，成功候选进入验证排序。每次运行建立新目录，记录配置、数据与代码哈希、依赖和成功/失败状态。生产式模拟通过后台线程运行事件循环，Web
请求经命令队列进入引擎；训练在子进程进行，检测继续使用旧模型。模型通过验证门槛并成功加载、预热后才替换内存中的模型引用。

\subsection{逻辑回归、Torch
线性模型与随机森林}\label{ux903bux8f91ux56deux5f52torch-ux7ebfux6027ux6a21ux578bux4e0eux968fux673aux68eeux6797}

逻辑回归和 Torch 线性模型都计算单个线性 logit，再经过 Sigmoid：

\[p(y=1\mid x)=\sigma(w^\mathsf{T}x+b),\qquad \sigma(z)=\frac{1}{1+e^{-z}}.\]

两者并不是同一次实验：历史逻辑回归使用 scikit-learn 的 \texttt{lbfgs}
求解器，族内获选参数为 \texttt{C=10}、不加权；Torch 线性模型是
\texttt{Linear(193,1)}，由 AdamW
和二元交叉熵训练。相同函数形式不意味着优化器、正则化和有限轮数下的结果相同。

随机森林由 200 棵决策树构成，对每棵树的类别概率求平均后以 0.5
判定。历史获选配置为不限制最大深度、叶节点最少 1 个样本，分裂特征使用
\texttt{sqrt}，采用 \texttt{balanced\_subsample}
类别权重。其结构为``预处理特征 → 多棵树并行判别 →
平均攻击概率''。树的实际深度由拟合数据决定，不将``不限深度''写成固定层数。多数类基线始终预测攻击，用于揭示单看
Recall 的不足。

\subsection{MLP 结构与训练}\label{mlp-ux7ed3ux6784ux4e0eux8badux7ec3}

正式原 MLP 和调参 MLP 都使用两层隐藏层，结构如下：

\[193\ \xrightarrow{\mathrm{Linear+ReLU}}\ 128
\ \xrightarrow{\mathrm{Linear+ReLU}}\ 64
\ \xrightarrow{\mathrm{Linear}}\ 1
\ \xrightarrow{\mathrm{Sigmoid}}\ p.\]

原 MLP 不使用 Dropout，训练 20 轮；调参获选版本在每个隐藏层 ReLU 后使用
Dropout 0.1，训练 40 轮。二者均使用学习率 0.001、权重衰减
0.0001、训练批量 512 和不加权 BCEWithLogitsLoss。推理时关闭
Dropout。数据标准化与类别编码保持一致，改进与退化均根据真实指标记录。

模拟服务另使用 MLP 64/32、CPU 2 线程、5 轮、训练批量
256。其输入维度取决于模拟训练池中的类别集合，不能直接套用完整训练集的
193 维；它也不是离线调参获选的 MLP 128/64，更不是 Transformer。

\subsection{轻量表格
Transformer}\label{ux8f7bux91cfux8868ux683c-transformer}

当前保留的是项目自定义轻量实现，源码见
\href{https://github.com/ZureTz/ml-network-intrusion-detection/blob/master/src/ml_network_intrusion_detection/training/transformer.py}{transformer.py}。它采用自注意力学习同一条记录内部的特征交互，参考
Transformer 与表格 Transformer 思路，但不是原始机器翻译
Transformer，也不是 FT-Transformer
论文的完整复现。本报告不包含已撤回论文版模型的训练或指标。

39 个标准化数值分别生成带字段身份的 token；154 维类别 One-Hot
向量联合投影成一个类别 token；加上一个可学习 CLS token，共 41 个
token。数值 token 和类别 token 可写为：

\[t_j=x_jw_j+b_j\in\mathbb{R}^{32},\qquad
 t_{\mathrm{cat}}=W_{\mathrm{cat}}x_{\mathrm{onehot}}+b_{\mathrm{cat}}.\]

\begin{figure}
\centering
\includegraphics[width=0.95\linewidth,height=\textheight,keepaspectratio,alt={轻量
Transformer
的实际数据维度与分类路径}]{final/figures/transformer-architecture.pdf}
\caption{轻量 Transformer 的实际数据维度与分类路径}
\end{figure}

族内获选网络宽度 32、4 个注意力头、2 个 Pre-Norm 块，单头维度
8。每个块先 LayerNorm，再生成
Q/K/V，执行缩放点积注意力，投影并加残差；随后 LayerNorm、32→128→32 的
GELU 前馈网络及残差。注意力概率、输出投影和前馈网络使用 Dropout
0.1。分类头取最后的 CLS，执行 LayerNorm 和 Linear(32,1)。

\[\operatorname{Attention}(Q,K,V)=\operatorname{softmax}\!\left(\frac{QK^\mathsf{T}}{\sqrt{d_k}}\right)V.\]

输入张量的形状为批量大小 × 41 × 32，注意力矩阵为每个头的 41 ×
41；批量中不同流量记录互不做注意力。训练 15 轮、学习率 0.0003、权重衰减
0.0001，网络批量上限 256。相对完整 FT-Transformer，当前类别字段合为单个
token，前馈激活和网络配置也不同，不能以相同名称暗示论文结果已复现。

\section{实验协议与候选选择}\label{ux5b9eux9a8cux534fux8baeux4e0eux5019ux9009ux9009ux62e9}

\subsection{环境与测量口径}\label{ux73afux5883ux4e0eux6d4bux91cfux53e3ux5f84}

正式实验运行于 Apple M1 Pro、8 个逻辑 CPU、16 GiB 内存、macOS
arm64；Python 3.12.13、PyTorch 2.15.0、scikit-learn
1.9.1。历史随机森林和逻辑回归使用 CPU，Torch 正式实验使用 MPS，配置 4 个
CPU 算子线程、float32。完整构建信息见各实验目录的
\texttt{environment.json}；CUDA 路径已实现但未在 NVIDIA 硬件上验证。

拟合耗时包含预处理、初始化和优化；GPU
计时在结束前同步。预测先预热一批，再对完整测试集预测三次，报告中位数。预测耗时包含
CPU 预处理、数据传输及同步后的模型计算，不包括模型加载、CSV
读写和特征提取；加载耗时单列。CSV 外层批量为 10,000，Transformer
内部网络批量上限 256，其他 Torch 神经模型为
10,000。因此性能反映当前配置，不能据此计算架构无关的 GPU 加速比。

内存每 50 ms 采样 RSS，具体进程/子进程可见性见原始记录；历史
scikit-learn 运行受限于当前进程采样。RSS
包含数据、解释器和分配器保留内存，且候选顺序执行，不是模型净内存或 GPU
显存。MPS 显存峰值没有测量。

\subsection{选择协议与实验批次}\label{ux9009ux62e9ux534fux8baeux4e0eux5b9eux9a8cux6279ux6b21}

所有正式实验固定种子 42 和阈值 0.5。只按验证集攻击 F1
降序、误报率升序、Recall
降序和候选顺序选优。选择冻结后，每个模型族只评估验证获胜候选；不合并训练/验证重拟合，不根据测试成绩替换
\texttt{best/}。

\Needspace{6\baselineskip}

{\def\LTcaptype{none} % do not increment counter
\begin{longtable}[]{@{}
  >{\raggedright\arraybackslash}p{(\linewidth - 6\tabcolsep) * \real{0.2500}}
  >{\raggedright\arraybackslash}p{(\linewidth - 6\tabcolsep) * \real{0.2500}}
  >{\raggedright\arraybackslash}p{(\linewidth - 6\tabcolsep) * \real{0.2500}}
  >{\raggedright\arraybackslash}p{(\linewidth - 6\tabcolsep) * \real{0.2500}}@{}}
\toprule\noalign{}
\begin{minipage}[b]{\linewidth}\raggedright
批次
\end{minipage} & \begin{minipage}[b]{\linewidth}\raggedright
候选范围
\end{minipage} & \begin{minipage}[b]{\linewidth}\raggedright
验证整体获选
\end{minipage} & \begin{minipage}[b]{\linewidth}\raggedright
证据
\end{minipage} \\
\midrule\noalign{}
\endhead
\bottomrule\noalign{}
\endlastfoot
2026-11-03 & 多数类 1、逻辑回归 6、森林 4 & 随机森林 &
\href{https://github.com/ZureTz/ml-network-intrusion-detection/blob/master/reports/results/full-20261103/comparison.csv}{历史结果} \\
2026-11-05 原 Torch & 多数类 1、线性 2、MLP 4 & MLP 128/64，无 Dropout &
\href{https://github.com/ZureTz/ml-network-intrusion-detection/blob/master/reports/results/torch-full-20261105/comparison.csv}{原
Torch} \\
2026-11-05 调参 & 多数类 1、线性 3、MLP 5、Transformer 2 & MLP
128/64，Dropout 0.1 &
\href{https://github.com/ZureTz/ml-network-intrusion-detection/blob/master/reports/results/torch-tune-20261105/comparison.csv}{调参结果} \\
\end{longtable}
}

历史逻辑回归比较 \texttt{C=\{0.1,1,10\}}
与是否加权；森林比较最大深度不限/20 和叶节点最少 1/5。原 Torch 比较 MLP
64/32、128/64，以及是否按训练类别计数加权。调参搜索在运行前固定，详见
\href{https://github.com/ZureTz/ml-network-intrusion-detection/blob/master/reports/results/torch-tune-20261105/training/search_plan.json}{search\_plan.json}。

调参线性比较学习率、权重衰减及 20/40 轮；MLP 比较
64/32、128/64、256/128、学习率、权重衰减、Dropout 和 20/40
轮；Transformer 比较 16 维/2 头/1 层和 32 维/4 头/2 层，均训练 15
轮。该搜索有限且预算不等，不代表各架构的全局最优比较。

\subsection{调参验证结果}\label{ux8c03ux53c2ux9a8cux8bc1ux7ed3ux679c}

下表为调参批次全部候选的真实验证记录；F1、Recall、FPR 均为百分数。

\Needspace{14\baselineskip}

{\def\LTcaptype{none} % do not increment counter
\begin{longtable}[]{@{}lrrrr@{}}
\toprule\noalign{}
候选 & 验证 F1 & 验证 Recall & 验证 FPR & 拟合秒数 \\
\midrule\noalign{}
\endhead
\bottomrule\noalign{}
\endlastfoot
dummy & 80.36 & 100.00 & 100.00 & 1.01 \\
linear-original & 95.15 & 98.87 & 18.31 & 10.09 \\
linear-lr003-e40 & 95.24 & 98.76 & 17.66 & 17.17 \\
linear-lr01-wd001-e40 & 95.29 & 99.04 & 18.06 & 17.17 \\
mlp-original & 95.90 & 95.96 & 8.54 & 14.79 \\
mlp128-lr0005-wd001-e40 & 95.88 & 98.58 & 14.41 & 32.36 \\
mlp128-drop01-e40 & 96.03 & 98.42 & 13.43 & 21.25 \\
mlp256-drop01-e40 & 95.97 & 98.16 & 13.10 & 27.11 \\
mlp64-lr003-e40 & 95.99 & 97.68 & 11.94 & 21.88 \\
transformer16-l1 & 95.71 & 96.92 & 11.48 & 114.61 \\
transformer32-l2 & 95.84 & 95.49 & 7.74 & 182.98 \\
\end{longtable}
}

本轮 11 个候选均成功，完整训练、验证与保存共 465.74 秒。整体获选
\texttt{mlp128-drop01-e40} 的验证 F1 为 96.03\%，高于 Transformer 的
95.84\%。训练时 Git 基线为
\nolinkurl{fe58aba02978ee7dc7bdb3757e192ad551991982}
且工作区有修改，精确代码依据运行记录中的源码
SHA-256，而不是把基线提交当成全部训练实现。原 Torch 批次的干净训练提交为
\nolinkurl{40ae60151fe68049b9a7a7ae2f754a37a169224a}。

\section{离线实验结果与分析}\label{ux79bbux7ebfux5b9eux9a8cux7ed3ux679cux4e0eux5206ux6790}

\subsection{指标定义}\label{ux6307ux6807ux5b9aux4e49}

以攻击为正类，混淆矩阵按真实类别为行、预测类别为列，顺序均为正常、攻击：

\[C=\begin{bmatrix}TN&FP\\FN&TP\end{bmatrix}.\]

\[\mathrm{Accuracy}=\frac{TP+TN}{N},\quad
\mathrm{Precision}=\frac{TP}{TP+FP},\quad
\mathrm{Recall}=\frac{TP}{TP+FN}.\]

\[F_1=\frac{2TP}{2TP+FP+FN},\qquad \mathrm{FPR}=\frac{FP}{FP+TN}.\]

宏平均 F1 为正常和攻击两类 F1 的等权平均。ROC-AUC 使用连续分数；Average
Precision（AP）按召回增量加权精度，定义为
\(\sum_n(R_n-R_{n-1})P_n\)，不等同于 PR
曲线梯形面积。零分母等未定义指标在 JSON 中保留 \texttt{null}
及原因，不填伪造数值。指标定义参见
\href{https://scikit-learn.org/stable/modules/generated/sklearn.metrics.average_precision_score.html}{scikit-learn
AP 文档}。

\subsection{各模型测试对比}\label{ux5404ux6a21ux578bux6d4bux8bd5ux5bf9ux6bd4}

以下指标均来自同一 82,332 行测试集，Accuracy、Precision、Recall、F1、FPR
均为百分数。合并表保留每一行的原始结果路径，见
\href{https://github.com/ZureTz/ml-network-intrusion-detection/blob/master/reports/final/comparison.csv}{comparison.csv}。

\Needspace{11\baselineskip}

{\def\LTcaptype{none} % do not increment counter
\begin{longtable}[]{@{}lrrrrr@{}}
\toprule\noalign{}
模型 & Accuracy & Precision & Recall & F1 & FPR \\
\midrule\noalign{}
\endhead
\bottomrule\noalign{}
\endlastfoot
多数类 & 55.06 & 55.06 & 100.00 & 71.02 & 100.00 \\
逻辑回归 & 81.13 & 75.52 & 97.25 & 85.02 & 38.61 \\
随机森林 & 86.83 & 81.31 & 98.78 & 89.20 & 27.82 \\
原 Torch 线性 & 80.23 & 74.70 & 96.90 & 84.36 & 40.20 \\
调参 Torch 线性 & 80.75 & 75.00 & 97.55 & 84.80 & 39.84 \\
原 MLP & 87.05 & 82.74 & 96.65 & 89.16 & 24.70 \\
调参 MLP & 84.45 & 78.55 & 98.71 & 87.49 & 33.02 \\
轻量 Transformer & 88.69 & 83.93 & 98.28 & 90.54 & 23.06 \\
\end{longtable}
}

\begin{figure}
\centering
\includegraphics[width=0.95\linewidth,height=\textheight,keepaspectratio,alt={同一测试集上的攻击
F1 与正常流量误报率}]{final/figures/model-comparison.pdf}
\caption{同一测试集上的攻击 F1 与正常流量误报率}
\end{figure}

多数类模型的攻击 Recall 为 100\%，但误报率同样为
100\%，不能作为有效检测器。轻量 Transformer 在已评估模型族中取得最高测试
F1，较原 MLP 高 1.39 个百分点，误报率低 1.64 个百分点；较随机森林 F1
高约 1.34 个百分点，但攻击 Recall
略低。以上是固定基准的描述性结果，不是统计显著性结论，也没有改变已冻结的验证选择。

\Needspace{11\baselineskip}

{\def\LTcaptype{none} % do not increment counter
\begin{longtable}[]{@{}lrrr@{}}
\toprule\noalign{}
模型 & 宏平均 F1（\%） & ROC-AUC & AP \\
\midrule\noalign{}
\endhead
\bottomrule\noalign{}
\endlastfoot
多数类 & 35.51 & 0.5000 & 0.5506 \\
逻辑回归 & 79.77 & 0.9561 & 0.9662 \\
随机森林 & 86.16 & 0.9794 & 0.9826 \\
原 Torch 线性 & 78.73 & 0.9478 & 0.9615 \\
调参 Torch 线性 & 79.27 & 0.9550 & 0.9667 \\
原 MLP & 86.55 & 0.9754 & 0.9819 \\
调参 MLP & 83.48 & 0.9765 & 0.9827 \\
轻量 Transformer & 88.24 & 0.9851 & 0.9888 \\
\end{longtable}
}

Transformer 的 ROC-AUC 和 AP 较高，但固定阈值下仍把约 23\%
的正常流量误判为攻击。不能以排序指标替代当前决策阈值下的误报分析。

\subsection{MLP
调参前后比较}\label{mlp-ux8c03ux53c2ux524dux540eux6bd4ux8f83}

\begin{figure}
\centering
\includegraphics[width=0.95\linewidth,height=\textheight,keepaspectratio,alt={MLP
调参前后的验证 F1、测试 F1 与测试误报率}]{final/figures/mlp-tuning.pdf}
\caption{MLP 调参前后的验证 F1、测试 F1 与测试误报率}
\end{figure}

原 MLP 验证 F1 为 95.8975\%，调参后为 96.0277\%，提高约 0.13
个百分点；测试 F1 从 89.1557\% 降为 87.4873\%，下降约 1.67
个百分点；误报率从 24.6973\% 升为 33.0162\%，上升约 8.32
个百分点。调参模型攻击 Recall 提高至 98.71\%，但产生了更多正常流量误报。

这种结果表明本轮调参未改善总体测试表现。单次实验不能区分过拟合、分布差异、有限训练预算或校准误差的因果贡献。保留原模型和真实结果比只报告获选参数更有助于解释泛化边界。测试集已被多轮实验使用，后续改进应在新协议和独立数据上确认。

\Needspace{12\baselineskip}

\subsection{混淆矩阵与攻击类别错误}\label{ux6df7ux6dc6ux77e9ux9635ux4e0eux653bux51fbux7c7bux522bux9519ux8bef}

\Needspace{7\baselineskip}

{\def\LTcaptype{none} % do not increment counter
\begin{longtable}[]{@{}lrrrr@{}}
\toprule\noalign{}
模型 & TN & FP & FN & TP \\
\midrule\noalign{}
\endhead
\bottomrule\noalign{}
\endlastfoot
随机森林 & 26706 & 10294 & 552 & 44780 \\
原 MLP & 27862 & 9138 & 1520 & 43812 \\
调参 MLP & 24784 & 12216 & 584 & 44748 \\
轻量 Transformer & 28469 & 8531 & 778 & 44554 \\
\end{longtable}
}

Transformer 相比原 MLP 少误报 607 条正常记录、少漏检 742
条攻击记录。其正常类 Recall 为 76.94\%，正常类 F1 为 85.95\%；攻击类
Recall 为
98.28\%。二分类错误细分如下，检出率指该攻击类别被判为攻击的比例，不表示能够识别其具体攻击名称。

\Needspace{12\baselineskip}

{\def\LTcaptype{none} % do not increment counter
\begin{longtable}[]{@{}lrrr@{}}
\toprule\noalign{}
攻击背景 & 测试条数 & 检出率（\%） & 漏报条数 \\
\midrule\noalign{}
\endhead
\bottomrule\noalign{}
\endlastfoot
Analysis & 677 & 99.85 & 1 \\
Backdoor & 583 & 99.83 & 1 \\
DoS & 4089 & 99.63 & 15 \\
Exploits & 11132 & 99.51 & 55 \\
Fuzzers & 6062 & 88.67 & 687 \\
Generic & 18871 & 99.98 & 3 \\
Reconnaissance & 3496 & 99.60 & 14 \\
Shellcode & 378 & 99.47 & 2 \\
Worms & 44 & 100.00 & 0 \\
\end{longtable}
}

Fuzzers 仍是主要漏报来源，6,062 条中漏报 687 条，占 Transformer 全部 778
条漏报的大部分。Worms 只有 44
条，即使全部检出也不能据此推断稳定泛化。原始测试冲突标签被保留，相同特征的确定性预测无法同时满足冲突标签。完整结果见
\href{https://github.com/ZureTz/ml-network-intrusion-detection/blob/master/reports/results/torch-tune-20261105/evaluation/transformer/metrics.json}{Transformer
指标}和\href{https://github.com/ZureTz/ml-network-intrusion-detection/blob/master/reports/results/torch-tune-20261105/evaluation/transformer/category_breakdown.csv}{类别细分}。

\subsection{训练、加载与推理成本}\label{ux8badux7ec3ux52a0ux8f7dux4e0eux63a8ux7406ux6210ux672c}

\Needspace{10\baselineskip}

{\def\LTcaptype{none} % do not increment counter
\begin{longtable}[]{@{}lrrrr@{}}
\toprule\noalign{}
模型 & 拟合（s） & 加载（s） & 预测（s） & 吞吐（行/s） \\
\midrule\noalign{}
\endhead
\bottomrule\noalign{}
\endlastfoot
逻辑回归 & 3.345 & 0.1069 & 0.1065 & 773095 \\
随机森林 & 28.679 & 0.3513 & 0.4954 & 166199 \\
原 Torch 线性 & 7.806 & 0.0250 & 0.1756 & 468902 \\
调参 Torch 线性 & 17.171 & 0.0237 & 0.1676 & 491289 \\
原 MLP & 10.856 & 0.0291 & 0.1865 & 441533 \\
调参 MLP & 21.248 & 0.0291 & 0.1815 & 453667 \\
轻量 Transformer & 182.983 & 0.0446 & 3.2356 & 25445 \\
\end{longtable}
}

表中训练及预测单位为秒，吞吐为行/秒；全部预测覆盖 82,332 行。原 MLP
的耗时取原 Torch
批次，其他模型取对应来源，不能把跨批次波动当成调参的确定收益。Transformer
更高的当前测试 F1 伴随更大计算成本，其完整测试预测约 3.236 秒，调参 MLP
约 0.181 秒。批处理吞吐不是单条网络事件时延，更不是生产链路容量。

独立 CPU 进程已经对原 MLP、调参 MLP、Transformer 的 82,332
行无标签输入验证了加载和预测：标签与行顺序一致，概率在
\texttt{atol=1e-6,\ rtol=1e-5} 内一致。验证对象是同一 CPU
加载方式下的直接预测与分批命令输出，不宣称 CPU 与 MPS 逐位相同。5
条训练未见的 \texttt{state} 已统计并正常处理。见
\href{https://github.com/ZureTz/ml-network-intrusion-detection/blob/master/reports/results/torch-tune-20261105/verification/transformer/verification.json}{Transformer
独立验证}。

\section{常驻推理、定时训练模拟与前端}\label{ux5e38ux9a7bux63a8ux7406ux5b9aux65f6ux8badux7ec3ux6a21ux62dfux4e0eux524dux7aef}

\subsection{模拟机制与运行方式}\label{ux6a21ux62dfux673aux5236ux4e0eux8fd0ux884cux65b9ux5f0f}

执行
\texttt{uv\ run\ ml-network-intrusion-detection\ serve\ -\/-config\ configs/simulation.yaml}
后启动本地常驻服务，访问终端给出的地址即可查看
Dashboard。模型加载后保留在进程内存中，对后续记录重复预测；关闭浏览器不停止服务，终端退出则保存状态。它模拟的是已提取特征的到达和标签反馈，而不是监听真实网络接口。

\begin{figure}
\centering
\includegraphics[width=0.95\linewidth,height=\textheight,keepaspectratio,alt={先预测、后揭示标签、再训练和验证的模拟流程}]{final/figures/simulation-flow.pdf}
\caption{先预测、后揭示标签、再训练和验证的模拟流程}
\end{figure}

初始化只读取已准备的 \texttt{train.csv} 和 \texttt{validation.csv}
并核验哈希。训练侧按标签分层抽取完整特征组形成约 10,000 条样本池，其中约
2,000 条用于起始模型；验证侧另取约 2,000
条固定样本。完整组约束使本次实际池大小为 10,007，起始样本为
2,060，固定验证集为 2,000。官方测试文件不进入模拟实验。

\Needspace{9\baselineskip}

{\def\LTcaptype{none} % do not increment counter
\begin{longtable}[]{@{}lll@{}}
\toprule\noalign{}
参数 & 默认/本次配置 & 含义 \\
\midrule\noalign{}
\endhead
\bottomrule\noalign{}
\endlastfoot
到达速率 & 初始 20 条/秒，可在 UI 调整 & 历史截图时为 150 条/秒 \\
标签延迟 & 10 个模拟秒 & 先预测，再揭示原始标签 \\
调度间隔 & 60 个模拟秒 & 检查是否可以启动新训练 \\
新样本门槛 & 500 条 & 新增且已确认的来源记录 \\
模型与训练 & MLP 64/32，5 轮，批量 256 & CPU，2 线程，从头拟合 \\
切换门槛 & 验证 F1 严格升高且 FPR 不增加 & 两项必须同时满足 \\
\end{longtable}
}

系统以来源 ID
去重，重复回放不会无限增加训练样本数。任务启动时冻结所有已揭示样本，形成累计训练快照；固定验证集不进入训练。子进程完成后比较新旧模型，在门槛成立、产物加载和预热成功后替换模型，否则保留旧版本。初始模型没有前任，首次成功训练直接装载。

这实现了``新数据不断到达，然后定期训练''的演示闭环，但有限样本池耗尽后只继续检测，来源去重使新样本数不再增加，不会无数据地反复训练。当前没有扰动字段或生成全新模拟样本的模块。原始比例、攻击升高与突发场景改变的是抽样比例和到达节奏，而非网络攻击生成机制。

\subsection{历史演示结果}\label{ux5386ux53f2ux6f14ux793aux7ed3ux679c}

本节使用 2026-11-06 演示留存的数据库与训练记录，2026-11-13
导出公开版证据。所有训练快照样本数已与 \texttt{source\_ids}
数量交叉核验，候选决策已与对应 \texttt{promotion.json} 核对。证据见
\href{https://github.com/ZureTz/ml-network-intrusion-detection/blob/master/reports/results/simulation-20261106/snapshot.json}{模拟快照}、\href{https://github.com/ZureTz/ml-network-intrusion-detection/blob/master/reports/results/simulation-20261106/experiment.json}{实验配置}和\href{https://github.com/ZureTz/ml-network-intrusion-detection/blob/master/reports/results/simulation-20261106/export_manifest.json}{来源哈希清单}。

\Needspace{11\baselineskip}

{\def\LTcaptype{none} % do not increment counter
\begin{longtable}[]{@{}lrrrrr@{}}
\toprule\noalign{}
任务 & 训练行数 & 旧 F1（\%） & 候选 F1（\%） & 候选 FPR（\%） & 决策 \\
\midrule\noalign{}
\endhead
\bottomrule\noalign{}
\endlastfoot
1 & 2060 & - & - & - & 初始装载 \\
2 & 3059 & 93.74 & 93.84 & 25.11 & 拒绝 \\
3 & 4260 & 93.74 & 94.67 & 22.68 & 接受 \\
4 & 5460 & 94.67 & 94.97 & 21.31 & 接受 \\
5 & 6661 & 94.97 & 95.16 & 19.79 & 接受 \\
6 & 7862 & 95.16 & 95.58 & 16.59 & 接受 \\
7 & 9062 & 95.58 & 95.28 & 16.89 & 拒绝 \\
8 & 10007 & 95.58 & 95.38 & 18.11 & 拒绝 \\
\end{longtable}
}

8 次任务包含首次训练、4 次后续接受和 3 次拒绝。3,059 条任务的候选 F1
虽略升，FPR 却从 24.35\% 升到 25.11\%，因此被拒绝；7,862 条任务同时改善
F1 和 FPR，被接受；9,062 和 10,007 条任务未通过门槛，系统继续使用 7,862
条时的版本。这证明验证门槛实际参与切换，而非训练完成就无条件替换。

快照累计检测事件 20,412 条，其中已获得反馈 18,907 条、待反馈 1,505
条；训练样本按来源去重后为 10,007 条。累计已反馈事件的 TN/FP/FN/TP 为
4,976/1,173/149/12,609，对应回放 Accuracy 93.01\%、F1 95.02\%、FPR
19.08\%。这些事件含重复回放，预测来自多个历史版本，不能与离线独立测试分数混在同一排名中。固定模拟验证集也被反复用于切换，因此其分数不是新的无偏泛化估计。

\subsection{可视化界面}\label{ux53efux89c6ux5316ux754cux9762}

\begin{figure}
\centering
\includegraphics[width=0.95\linewidth,height=\textheight,keepaspectratio,alt={历史
Dashboard
总览：事件计数、标签反馈、训练日志和当前模型}]{final/figures/dashboard-overview.png}
\caption{历史 Dashboard 总览：事件计数、标签反馈、训练日志和当前模型}
\end{figure}

前端提供暂停/恢复、速率与场景调整、手动触发有新样本的训练、独立实验创建、攻击告警过滤、模型版本筛选和摘要导出。总览中的吞吐由浏览器观察事件数估算；截图时处于暂停展示状态，不能把图中瞬时速率当成系统峰值。

\begin{figure}
\centering
\includegraphics[width=0.95\linewidth,height=\textheight,keepaspectratio,alt={事件与告警列表：标签尚未揭示时显示待确认}]{final/figures/dashboard-events.png}
\caption{事件与告警列表：标签尚未揭示时显示待确认}
\end{figure}

事件列表展示来源行、预测结果、攻击分数、真实标签是否已反馈、模型版本与整批预测耗时。截图中的
3.5 ms
属于一批模型预测时间，并非每条记录的端到端网络延迟。表中没有姓名、账户或原始
IP。

\begin{figure}
\centering
\includegraphics[width=0.95\linewidth,height=\textheight,keepaspectratio,alt={模型验证与版本记录：接受与拒绝均可追溯}]{final/figures/dashboard-jobs.png}
\caption{模型验证与版本记录：接受与拒绝均可追溯}
\end{figure}

版本记录保留训练样本规模和新旧验证指标，有助于解释为何更新或保留模型。前端展示的训练日志来自实际子进程；本报告使用原始界面截图，不用人工绘制的页面代替运行证据。

\subsection{状态、日志与失败处理}\label{ux72b6ux6001ux65e5ux5fd7ux4e0eux5931ux8d25ux5904ux7406}

\texttt{state.sqlite}
保存模拟时钟、随机状态、来源去重、待反馈事件、训练历史和当前模型路径；每个训练目录保存快照、配置、日志、模型及决策。\texttt{service.log}
采用 10 MB 轮转、5 份备份。事件明细最多保留最近 100,000
条已确认记录，待反馈事件不清理；累计混淆计数独立保存。训练目录不自动删除，长期使用需要归档。

暂停会冻结模拟时钟及新任务调度，已启动任务可完成。退出会停止子进程，下次恢复将未完成任务标为中断。训练失败保留旧模型，30
秒后允许重试；单次训练最长 2 小时。待生成事件额度最多积压
1,000，超出部分记录为
dropped，这不是捕获网络包的丢包统计。固定种子可控制抽样，但系统调度会影响批次大小，不保证逐事件重现。

Web
服务只监听回环地址并限制同源控制请求，定位为单机实验工具。未验证远程多用户、生产级认证、真实标签反馈、特征抽取一致性或实际概念漂移适应性。

\section{工程验证、复现与公开交付}\label{ux5de5ux7a0bux9a8cux8bc1ux590dux73b0ux4e0eux516cux5f00ux4ea4ux4ed8}

\subsection{已完成的验证}\label{ux5df2ux5b8cux6210ux7684ux9a8cux8bc1}

2026-11-13 撤回论文版新增代码后，\texttt{uv\ sync\ -\/-locked}
成功，Ruff 检查与格式检查通过，完整测试为 54 passed、4
skipped。四个跳过项为当前沙箱不可用的 MPS/CUDA
硬件测试；不能据此判断主机从未支持 MPS。2026-11-05 正式实验及当时真实
MPS 测试已留有记录。CPU 环境检查完成了训练、保存/加载及概率一致性，旧
Transformer 正式权重又完成了 8 行无标签特征预测检查。

自动化测试覆盖数据组隔离、预处理不泄漏、标签映射、未知类别、缺少特征、保存加载、独立进程预测、注意力实际梯度、批量隔离、延迟标签、训练快照、模型切换、暂停与重启恢复。合成样本测试不作为检测能力证据。报告更新不重新训练模型，表格和图表从已保存结果生成；本次文档构建另检查链接、数据一致性、LaTeX
编译和逐页排版。

\subsection{复现命令与兼容性}\label{ux590dux73b0ux547dux4ee4ux4e0eux517cux5bb9ux6027}

在项目根目录恢复依赖，并先按
\href{https://github.com/ZureTz/ml-network-intrusion-detection/blob/master/README.md}{根目录说明}下载数据。以下目录名为新运行示例，已有输出不可覆盖。

\begin{verbatim}
uv sync --locked
uv run ml-network-intrusion-detection prepare \
  --config configs/full.yaml --output data/processed/report-run
uv run ml-network-intrusion-detection train \
  --config configs/tune.yaml --data data/processed/report-run \
  --output models/report-run --device mps
uv run python scripts/evaluate_selected.py \
  --training models/report-run --data data/processed/report-run \
  --output outputs/report-run --device mps
uv run python scripts/verify_inference.py \
  --model models/report-run/transformer32-l2 \
  --input data/processed/report-run/test.csv \
  --output outputs/report-run-verify --device cpu
\end{verbatim}

无 MPS 的环境可改用 CPU，但训练耗时和浮点结果可能不同。模拟启动前，应将
\texttt{configs/simulation.yaml} 中 \texttt{data}
改为新准备目录、\texttt{output}
改为新的实验目录，路径相对于配置文件；然后执行
\texttt{serve}。已有模拟目录恢复时沿用其保存的实验参数，不能仅修改 YAML
就假定旧实验已被重建。

当前 Torch 模型包含 \texttt{preprocess.joblib}、CPU \texttt{weights.pt}
与
\texttt{metadata.json}；加载时校验字段、类别、阈值、依赖和文件哈希，Torch
权重以 \texttt{weights\_only=True} 读取。joblib/pickle
仍只可加载可信来源。历史 scikit-learn 模型使用旧格式，不能直接放入当前
Torch 加载流程；复现历史实验需其对应历史代码和锁文件。

\subsection{报告构建和证据索引}\label{ux62a5ux544aux6784ux5efaux548cux8bc1ux636eux7d22ux5f15}

安装 Pandoc 与包含 XeLaTeX、ctex、Fandol 字体的 TeX Live 后：

\begin{verbatim}
uv run python scripts/plot_report.py
uv run python scripts/build_report.py
\end{verbatim}

图表脚本只读取已提交的精简指标，生成合并 CSV 和图表；构建脚本由本
Markdown 生成完整 LaTeX，再运行三遍
XeLaTeX。正文、表格、公式均为可编辑文本，LaTeX 与 PDF
一并提交。编译中间文件写入被忽略的
\texttt{outputs/report-build/}。报告统一使用常规字重，不使用加粗强调。

\Needspace{9\baselineskip}

{\def\LTcaptype{none} % do not increment counter
\begin{longtable}[]{@{}
  >{\raggedright\arraybackslash}p{(\linewidth - 2\tabcolsep) * \real{0.5000}}
  >{\raggedright\arraybackslash}p{(\linewidth - 2\tabcolsep) * \real{0.5000}}@{}}
\toprule\noalign{}
\begin{minipage}[b]{\linewidth}\raggedright
材料
\end{minipage} & \begin{minipage}[b]{\linewidth}\raggedright
用途
\end{minipage} \\
\midrule\noalign{}
\endhead
\bottomrule\noalign{}
\endlastfoot
\href{https://github.com/ZureTz/ml-network-intrusion-detection/blob/master/reports/sklearn-20261103.md}{历史
scikit-learn 报告} & 历史模型、配置和运行边界 \\
\href{https://github.com/ZureTz/ml-network-intrusion-detection/blob/master/reports/torch-20261105.md}{原
Torch 报告} & 原 MLP 与线性正式结果归档 \\
\href{https://github.com/ZureTz/ml-network-intrusion-detection/blob/master/reports/transformer-tuning-20261105.md}{调参专项报告}
& 完整调参过程与独立预测验证 \\
\href{https://github.com/ZureTz/ml-network-intrusion-detection/blob/master/reports/final/comparison.csv}{最终合并
CSV} & 模型族比较及每行结果来源 \\
\href{https://github.com/ZureTz/ml-network-intrusion-detection/blob/master/reports/results/simulation-20261106/}{模拟证据目录}
& 去除本机绝对路径的实验快照与哈希 \\
\href{https://github.com/ZureTz/ml-network-intrusion-detection/blob/master/reports/final/README.md}{构建说明}
& LaTeX/PDF 重建方法与检查口径 \\
\end{longtable}
}

公开交付只包括报告、LaTeX、PDF、必要图表及精简证据，不包括原始数据、模型二进制、逐行预测或私人
PPT。\texttt{reports/presentations/} 与 PPT 扩展名已加入 Git
忽略规则；报告截图已检查可见内容，不含个人姓名。已被 Git
跟踪的文件不会因忽略规则自动移除，因此推送前还需检查暂存清单。

\section{结论与局限}\label{ux7ed3ux8bbaux4e0eux5c40ux9650}

项目已完成可复现的离线入侵检测流程，并通过本地常驻服务展示特征到达、模型推理、延迟反馈、累计样本重训、验证和模型切换。模型结构、离线指标、时间成本及前端行为均有对应代码或运行证据。轻量
Transformer 在当前测试比较中具有较高 F1，原 MLP 则有较低推理成本；调参
MLP 未带来测试收益，保留这一结果有助于避免仅凭验证提升作结论。

主要局限是单数据集、单种子、有限且不等的搜索预算，缺少多种子置信区间、阈值校准、外部数据、主机/时间隔离和真实网络验证。测试集已被历史方案反复使用，结果是探索性基准比较。即使当前较好的
Transformer，正常流量误报率仍为 23.06\%，不能据此宣布满足生产告警要求。

模拟部分没有新独立数据生成能力、真实标签反馈来源或实际概念漂移实验。固定验证集的重复使用可能使切换决策逐渐适应该验证集；样本池耗尽后不再产生新训练样本。下一步应先定义独立数据和误报约束，再用验证集选择阈值、执行多种子对照并验证新的测试集，而非继续根据现有测试分数反复调整模型。

\clearpage

\section{参考文献}\label{ux53c2ux8003ux6587ux732e}

以下资料用于数据来源、架构背景和指标定义。官方数据页要求引用的相关文献一并保留；背景文献不代表本项目完整复现其全部实验。

\begin{enumerate}
\def\labelenumi{\arabic{enumi}.}
\tightlist
\item
  Moustafa, N.; Slay, J. UNSW-NB15: a comprehensive data set for network
  intrusion detection systems (UNSW-NB15 network data set). MilCIS,
  2015.
  \href{https://ieeexplore.ieee.org/abstract/document/7348942}{IEEE
  记录}。
\item
  Moustafa, N.; Slay, J. The evaluation of Network Anomaly Detection
  Systems: Statistical analysis of the UNSW-NB15 data set and the
  comparison with the KDD99 data set. Information Security Journal,
  2016. \href{https://doi.org/10.1080/19393555.2015.1125974}{DOI}。
\item
  Moustafa, N., et al.~Novel Geometric Area Analysis Technique for
  Anomaly Detection Using Trapezoidal Area Estimation on Large-Scale
  Networks. IEEE Transactions on Big Data，官方页列为 2017
  年在线资料。\href{https://ieeexplore.ieee.org/abstract/document/7948715}{IEEE
  记录}。
\item
  Moustafa, N.; Creech, G.; Slay, J. Big Data Analytics for Intrusion
  Detection System: Statistical Decision-Making Using Finite Dirichlet
  Mixture Models. 2017, pp.~127-156.
  \href{https://link.springer.com/chapter/10.1007/978-3-319-59439-2_5}{Springer
  章节}。
\item
  Sarhan, M.; Layeghy, S.; Moustafa, N.; Portmann, M. NetFlow Datasets
  for Machine Learning-based Network Intrusion Detection Systems.
  \href{https://arxiv.org/abs/2011.09144}{arXiv:2011.09144}。本项目未改用
  NetFlow 版本。
\item
  Vaswani, A., et al.~Attention Is All You Need. NeurIPS, 2017.
  \href{https://arxiv.org/abs/1706.03762}{arXiv:1706.03762}。
\item
  Gorishniy, Y.; Rubachev, I.; Khrulkov, V.; Babenko, A. Revisiting Deep
  Learning Models for Tabular Data. NeurIPS, 2021.
  \href{https://arxiv.org/abs/2106.11959}{arXiv:2106.11959}。当前模型仅借鉴表格注意力思路。
\item
  UNSW Canberra.
  \href{https://research.unsw.edu.au/projects/unsw-nb15-dataset}{The
  UNSW-NB15 Dataset}。数据与引用信息核对日期：2026-11-13。
\item
  scikit-learn.
  \href{https://scikit-learn.org/stable/modules/generated/sklearn.metrics.average_precision_score.html}{Average
  Precision 定义}。核对日期：2026-11-13。
\item
  PyTorch.
  \href{https://docs.pytorch.org/docs/2.15/generated/torch.nn.functional.scaled_dot_product_attention.html}{Scaled
  dot product
  attention}、\href{https://docs.pytorch.org/docs/2.15/notes/mps.html}{MPS
  backend}、\href{https://docs.pytorch.org/docs/2.15/notes/serialization.html}{Serialization
  semantics}。实现接口与持久化说明。
\end{enumerate}
\end{document}
